\documentclass[11pt]{article}

\usepackage[margin=1in]{geometry}
\usepackage{amsmath, amssymb, amsthm}
\usepackage{graphicx}
\usepackage{hyperref}

\newtheorem{theorem}{Theorem}
\newtheorem{corollary}{Corollary}

\title{An Irreducible Accuracy Ceiling Under Two-Class Swap Corruption}
\author{Cate Merfeld Dunham}
\date{\today}

\begin{document}
\maketitle

\subsection*{Overview}

This note grew out of an experiment where an encoder was trained to map MNIST
images of two digits (6 and 8) onto 10-dimensional one-hot targets, with 20\%
of the training labels deliberately swapped to the other digit. On held-out
test images -- never seen during training, and clean of label corruption --
the trained encoder never exceeded about 80--85\% accuracy against the true
digit label, no matter how long it was trained. That raised a natural
question: is this just a property of this particular encoder and training
run, or is 80\% a hard ceiling that \emph{no} classifier could beat at test
time, given training data corrupted this way? (Training-set accuracy is a
different quantity, discussed separately below, since a model can simply
memorize individual training examples.)

It turns out to be the latter, and the argument generalizes well beyond
MNIST. Picture a sensor distinguishing two chemicals, A and B, that is bad
enough that a fixed fraction $p$ of the time, a reading from chemical A comes
out looking \emph{exactly} like a normal reading of chemical B (and vice
versa) --- not just noisy, but a full masquerade, indistinguishable from the
genuine article. The claim is that in this situation, no classifier, however
clever, can be more than $1-p$ accurate at recovering the true chemical from
the reading alone. This document states that claim precisely as a theorem
and proves it, then connects the proof back to the empirical results.

\subsection{Setup and assumptions}
\label{sec:setup}

Let $Y \in \{A, B\}$ denote the true class (chemical, digit, etc.), with
equal priors $P(Y=A) = P(Y=B) = \tfrac12$. Let $X$ denote the observed
reading. The model has two ingredients:

\begin{itemize}
  \item \textbf{Clean signal.} Each class has a ``clean'' reading
  distribution, $f_A$ and $f_B$, with \emph{disjoint support} $S_A \cap S_B =
  \emptyset$. In other words, when the sensor is not fooled, the reading is a
  perfect fingerprint of the true class.
  \item \textbf{Swap corruption.} Independently of everything else, with
  probability $p \in [0, \tfrac12]$ the sensor is fooled: a sample from class
  $A$ produces a reading drawn from $f_B$ instead of $f_A$, and a sample from
  class $B$ produces a reading drawn from $f_A$ instead of $f_B$.
\end{itemize}

Formally,
\[
  X \mid Y = A \;\sim\; (1-p)\,f_A + p\,f_B,
  \qquad
  X \mid Y = B \;\sim\; (1-p)\,f_B + p\,f_A.
\]

This is the ``symmetric class-flip'' noise model studied in the
noisy-labels literature (e.g.\ Angluin and Laird, 1988; Natarajan et al.,
2013), applied here to a physical measurement rather than a labeling
process. The critical modeling choice is the \emph{swap}: a corrupted
reading is not merely uninformative, it is statistically identical to a
genuine reading of the other class. This is what turns the bound into an
exact equality rather than a looser inequality (Section~\ref{sec:discussion}
discusses the alternative, softer model).

\subsection{Main theorem}

\begin{theorem}[Irreducible error under two-class swap corruption]
\label{thm:main}
Under the model of Section~\ref{sec:setup}, for any $p \in [0, \tfrac12]$,
every classifier $g : \mathcal{X} \to \{A, B\}$ satisfies
\[
  P\big(g(X) = Y\big) \;\le\; 1 - p,
\]
and this bound is achieved by the classifier
\[
  g^\ast(x) = \begin{cases} A & x \in S_A \\ B & x \in S_B. \end{cases}
\]
\end{theorem}

In words: no classifier --- regardless of model class, capacity, or training
procedure --- can recover the true class more than $(1-p) \times 100\%$ of
the time, and a plug-in classifier which simply reports whichever class's
clean signature matches the observed reading already attains this ceiling.

\begin{quote}
\textbf{Scope: this is a test-time (generalization) bound.} The probability
$P(g(X)=Y)$ in Theorem~\ref{thm:main} is a population quantity: it is the
accuracy of $g$ on a \emph{fresh} draw of $(X,Y)$ from the corruption model,
not the accuracy $g$ achieves on whatever finite sample was used to fit it.
The two can differ. A classifier with enough capacity can memorize the
identity of each individual training example and simply look up its known
true label, which trivially reproduces $100\%$ accuracy on the training set
regardless of $p$ --- this is not a violation of the theorem, because a
lookup keyed on ``which specific training example is this'' uses
information beyond the reading's distribution, which the theorem's argument
assumes is unavailable. On a genuinely new reading (i.e.\ at test time, or
for any sample not seen during fitting), that side channel does not exist,
and the $1-p$ ceiling applies in full force. Section~\ref{sec:empirical}
therefore checks the theorem's prediction against \emph{test}-set accuracy
specifically, not training-set accuracy.
\end{quote}

\subsection{Proof}

\subsubsection*{Step 1: Bayes-optimality of the MAP rule}

Let $\eta(x) = P(Y = A \mid X = x)$ be the posterior probability of class
$A$. For any classifier $g$,
\[
  P(g(X) = Y)
  = \mathbb{E}_X\Big[\, \eta(X)\,\mathbf{1}\{g(X) = A\}
    + \big(1-\eta(X)\big)\,\mathbf{1}\{g(X) = B\} \,\Big]
  \;\le\; \mathbb{E}_X\big[\, \max\big(\eta(X),\, 1-\eta(X)\big) \,\big],
\]
since at each $x$ the best achievable value of the bracketed term is
$\max(\eta(x), 1-\eta(x))$, attained by predicting whichever class is more
likely at that $x$. This upper bound is a general fact about 0--1 loss
classification and does not yet use the corruption model; it is achieved
with equality by the Bayes (maximum a posteriori) classifier
$g^\ast(x) = \arg\max\big(\eta(x), 1-\eta(x)\big)$.

\subsubsection*{Step 2: the posterior under swap corruption}

Fix $x \in S_A$. Since $S_A \cap S_B = \emptyset$, we have $f_B(x) = 0$, so
\[
  P(X = x \mid Y = A) \propto (1-p)\,f_A(x),
  \qquad
  P(X = x \mid Y = B) \propto p\,f_A(x).
\]
With equal priors, Bayes' rule gives
\[
  \eta(x) = \frac{(1-p)f_A(x)}{(1-p)f_A(x) + p\,f_A(x)} = 1 - p,
  \qquad \text{for all } x \in S_A.
\]
By the symmetric argument, $\eta(x) = p$ for all $x \in S_B$. (Every reading
lands in $S_A$ or $S_B$: even a masquerading class-$A$ sample's reading is
literally drawn from $f_B$, hence lies in $S_B$.)

\subsubsection*{Step 3: combining}

For $p \le \tfrac12$, $\max(\eta(x), 1-\eta(x)) = 1-p$ on both $S_A$ and
$S_B$. Substituting into Step 1,
\[
  P(g(X) = Y) \;\le\; \mathbb{E}_X[\,1-p\,] \;=\; 1-p
\]
for every classifier $g$.

\subsubsection*{Step 4: achievability}

The classifier $g^\ast$ from Theorem~\ref{thm:main} is exactly the Bayes
rule from Step 1, since $\eta(x) > 1-\eta(x)$ on $S_A$ (as $1-p > p$) and
symmetrically on $S_B$. Its accuracy is
\[
  P(X \in S_A, Y=A) + P(X \in S_B, Y=B)
  = \tfrac12(1-p) + \tfrac12(1-p) = 1-p,
\]
matching the upper bound exactly. $\blacksquare$

\begin{corollary}
\label{cor:twenty}
At $p = 0.2$, no classifier can exceed $80\%$ accuracy, and $80\%$ is
achievable.
\end{corollary}

\subsection{Why the decision boundary cannot help}

The proof clarifies why moving the decision boundary within the ``smudge''
of ambiguous readings cannot improve accuracy: correctly catching one more
masquerading reading by biasing the boundary toward the true class forces a
one-for-one misclassification of a genuine reading of the other class that
shares the same appearance. Inside the ambiguous region the two classes are,
by construction, statistically identical, so no boundary placement has
leverage. This is the precise sense in which the $80\%$ ceiling is
\emph{hard}: it is not a property of a particular classifier or training
procedure, but a limit on the information present in the reading itself.

\subsection{Connection to the empirical results}
\label{sec:empirical}

The MNIST 6-vs-8 encoder experiment corrupts $20\%$ of training labels by
assigning them the \emph{other} digit's one-hot target -- exactly the swap
model of Section~\ref{sec:setup}, with the one-hot target playing the role
of the reading $X$ and the true digit playing the role of $Y$.
Theorem~\ref{thm:main} and Corollary~\ref{cor:twenty} predict a ceiling of
$80\%$ accuracy at test time. As noted above, this is specifically a
statement about \textbf{test-set} accuracy: the test images are never seen
during training and their labels are never corrupted, so the encoder cannot
have memorized anything sample-specific about them, and its output on them
must depend only on the (possibly ambiguous) information genuinely present
in the reading -- exactly the regime the theorem describes. Training-set
accuracy against the true label is a different, unbounded quantity, since
the encoder's parameters were fit directly on those same samples and could
in principle memorize each one's true digit individually, regardless of
what corrupted target its loss term used.

Figure~\ref{fig:pca} shows the trained encoder's outputs, reduced to two
dimensions via PCA, for both splits; the test panel (right) is the one the
theorem constrains. On the test split, the trained encoder reaches $84.13\%$
accuracy on digit 6 and $84.80\%$ on digit 8, for an overall test accuracy
of $84.47\%$ -- consistent with (and slightly above) the $80\%$ ceiling from
Corollary~\ref{cor:twenty}. For reference, training-set accuracy is
$81.39\%$ on digit 6 and $81.66\%$ on digit 8 ($81.53\%$ overall); as
discussed above, this quantity is not constrained by Theorem~\ref{thm:main}.

\begin{figure}[h]
  \centering
  \includegraphics[width=\textwidth]{pca_plots.png}
  \caption{PCA projection of encoder outputs against one-hot targets, for
  the training split (left, $20\%$ label corruption, not bounded by
  Theorem~\ref{thm:main}) and test split (right, clean labels, the split
  the theorem constrains). Per-digit classification accuracy is annotated
  above each panel.}
  \label{fig:pca}
\end{figure}

The measured \emph{test} accuracy sits at or slightly above the $80\%$
ceiling predicted by Corollary~\ref{cor:twenty}. The small excess is
expected: the disjoint-support assumption in Section~\ref{sec:setup} is an
idealization, since real digit images belonging to the clean $80\%$ are not
perfectly separable either, leaving a small amount of headroom above the
ceiling from genuine (small) class overlap, together with ordinary
finite-sample fluctuation.

\subsection{The smudge on held-out data}
\label{sec:discussion}

Theorem~\ref{thm:main} bounds the accuracy of a \emph{decision}, a choice
between two classes. The continuous gradient of points between the two target
clusters in Figure~\ref{fig:pca} concerns a different object: the
continuous-valued output of an encoder trained to minimize squared error
against one-hot targets. A second result addresses that object.

Let $X$ be the reading, taking values in a space $\mathcal{X}$, jointly
distributed with a random target $T$.

\begin{theorem}[Optimal squared-error prediction]
\label{thm:l2}
Let $T$ be a random vector taking one of two values, $e_A$ or $e_B$
($e_A, e_B \in \mathbb{R}^n$), with $P(T=e_A \mid X=x)=p(x)$ and
$P(T=e_B \mid X=x) = 1-p(x)$. Then, over all measurable
$f: \mathcal{X} \to \mathbb{R}^n$, the expected squared distance between
$f(X)$ and $T$, $\mathbb{E}\big[\|f(X) - T\|^2\big]$, is minimized (a.e.\
uniquely) by:
\begin{equation}
    f^*(x) = \mathbb{E}[T\mid X=x] = p(x)\,e_A + (1-p(x))\,e_B.
\end{equation}
\end{theorem}

\begin{proof}
\textbf{Step 1: a single reading.} Fix a reading $x$, and let
$v \in \mathbb{R}^n$ be a candidate output at that reading, that is, any
vector a predictor could output there. Given $X=x$, the target equals $e_A$ with probability $p(x)$ and $e_B$ with
probability $1-p(x)$, so
\[
  \mathbb{E}\big[\|v-T\|^2 \mid X=x\big]
  = p(x)\,\|v-e_A\|^2 + (1-p(x))\,\|v-e_B\|^2 .
\]
Expanding each term with $\|v-e\|^2 = \|v\|^2 - 2\,v\cdot e + \|e\|^2$ gives
\[
  \|v\|^2 - 2\,v\cdot f^*(x) + p(x)\|e_A\|^2 + (1-p(x))\|e_B\|^2 .
\]
Since $\|v-f^*(x)\|^2 = \|v\|^2 - 2\,v\cdot f^*(x) + \|f^*(x)\|^2$, this equals
\[
  \mathbb{E}\big[\|v-T\|^2 \mid X=x\big] = \|v-f^*(x)\|^2 + c(x),
\]
where
\[
  c(x) := p(x)\|e_A\|^2 + (1-p(x))\|e_B\|^2 - \|f^*(x)\|^2
        = p(x)\,(1-p(x))\,\|e_A-e_B\|^2
\]
does not depend on $v$. In particular, at the reading $x$ the output
$v=f^*(x)$ is the unique minimizer of the expected squared error, and the
smallest value attained is $c(x)$.

\textbf{Step 2: all readings.} Let $f$ be measurable. Since $f(X)$ is a
function of $X$, Step~1 applies with $v = f(X)$, giving
\[
  \mathbb{E}\big[\|f(X)-T\|^2 \mid X\big] = \|f(X)-f^*(X)\|^2 + c(X).
\]
Taking expectations and using the law of total expectation,
\[
  \mathbb{E}\big[\|f(X)-T\|^2\big]
  = \mathbb{E}\big[\|f(X)-f^*(X)\|^2\big] + \mathbb{E}\big[c(X)\big].
\]
Applying this to $f=f^*$ shows $\mathbb{E}[c(X)] = \mathbb{E}\big[\|f^*(X)-T\|^2\big]$,
which is finite because $c(x) \le \tfrac14\|e_A-e_B\|^2$. Hence
\begin{equation}
  \mathbb{E}\big[\|f(X)-T\|^2\big]
  = \mathbb{E}\big[\|f^*(X)-T\|^2\big] + \mathbb{E}\big[\|f(X)-f^*(X)\|^2\big],
  \label{eq:decomp}
\end{equation}
as an identity in $[0,\infty]$. Because
$\mathbb{E}\big[\|f(X)-f^*(X)\|^2\big] \ge 0$, no $f$ can do better than
$f^*$. If $f$ attains the minimum, then
$\mathbb{E}\big[\|f(X)-f^*(X)\|^2\big] = 0$, and a nonnegative random variable
with zero mean is zero almost surely, so $f(X)=f^*(X)$ almost surely.
\end{proof}

\begin{corollary}[The smudge]
\label{cor:smudge}
At every reading $x$ with $0<p(x)<1$, the optimal output $f^*(x)$ lies
strictly between $e_A$ and $e_B$ on the segment joining them. Any measurable
$f$ with $P\big(f(X)\neq f^*(X)\big) > 0$, in particular any $f$ that outputs
$e_A$ or $e_B$ on a set of readings of positive probability where
$0<p(x)<1$, satisfies
\[
  \mathbb{E}\big[\|f(X)-T\|^2\big] > \mathbb{E}\big[\|f^*(X)-T\|^2\big].
\]
\end{corollary}

\begin{proof}
For $0<p(x)<1$, $f^*(x)$ is a convex combination of $e_A$ and $e_B$ with both
weights strictly positive, hence lies in the open segment and equals neither
endpoint. If $f(X)\neq f^*(X)$ with positive probability, then
$\|f(X)-f^*(X)\|^2 > 0$ on that event, so
$\mathbb{E}\big[\|f(X)-f^*(X)\|^2\big] > 0$, and the strict inequality follows
from~\eqref{eq:decomp}.
\end{proof}

Both statements are expectations over the data distribution, so they describe
error on \emph{new} samples, not on the finite sample used for fitting. This
matters for interpretation. On the training set, a network with enough
capacity can memorize each sample's target and avoid any smudge; that does
not contradict the result, because a training-set average is not the
population expectation. On held-out data, the smudge is not an artifact of
imperfect training: wherever the readings are genuinely ambiguous, it is what
optimal performance requires, and a model that avoids it is measurably worse.

Two limits on this claim should be kept in view. The result does not say that
every trained network smudges; it says smudging is required to be optimal, and
that departures from it cost a computable amount. And it forces a smudge only
at readings whose class is genuinely uncertain under the held-out
distribution. When the held-out labels are clean, as in the test set used
here, ambiguity comes only from images that are truly hard to tell apart, not
from the label corruption applied during training. The smudge is unrelated to
the value of the accuracy ceiling in Theorem~\ref{thm:main}, which is
unchanged.

\subsection{Extensions}

Two natural generalizations of Theorem~\ref{thm:main} are worth stating for
future work:

\begin{itemize}
  \item \textbf{Unequal priors.} If $P(Y=A) = \pi \neq \tfrac12$, the
  posterior on $S_A$ becomes
  $\eta(x) = \dfrac{\pi(1-p)}{\pi(1-p) + (1-\pi)p}$,
  which no longer simplifies to $1-p$; the ceiling then depends on both
  $p$ and $\pi$.
  \item \textbf{Asymmetric corruption rates.} If class $A$ is corrupted at
  rate $p_A$ and class $B$ at a different rate $p_B$, the two branches of
  the proof no longer share the same value of $\max(\eta, 1-\eta)$, and the
  overall ceiling becomes a weighted average of $(1-p_A)$ and $(1-p_B)$
  rather than a single shared constant.
  \item \textbf{Soft (non-swap) ambiguity.} If corrupted readings are drawn
  from a shared, genuinely uninformative distribution instead of exactly
  mimicking the other class, the ceiling relaxes to $1 - p/2$: the clean
  $(1-p)$ fraction is still recovered exactly, while the ambiguous $p$
  fraction can only be guessed at chance level.
\end{itemize}

\end{document}
